% ─── Sensitivity Analysis ────────────────────────────────────────────────────
\label{sec:sensitivity}

The central findings rest on a set of uncertain parameters and structural assumptions,
and it is necessary to establish which of them the conclusions actually turn on. This
section examines that sensitivity through six methods, hierarchically organised and
grouped by the side of the model they probe. The first four address the demand side: a
one-at-a-time tornado on the three economic axes that govern the technology mix, the
ETS2 carbon-price trajectory, the green-hydrogen import price and the discount rate,
reported for 2030 buildings-sector emissions under Stated Policies
(Appendix~\ref{sec:sensitivity:tornado}); a parametric band on the cross-country
renovation correlation $\rho$, which sets the width of the demand distribution
(Appendix~\ref{sec:sensitivity:rho}); a variance-based (Sobol) global decomposition that
apportions the variance of 2050 useful heat across its four demand-side drivers
(Appendix~\ref{sec:sensitivity:sobol}); and a one-sided climate-warming bound on demand
(Appendix~\ref{sec:sensitivity:climate}).

The remaining two are targeted stress tests on the supply side: on the per-country
hydrogen-supply cost (Appendix~\ref{sec:sensitivity:h2supply}) and on the electricity
network reinforcement and scarcity premium (Appendix~\ref{sec:sensitivity:reinforce}).
The methods are complementary. The tornado and the parametric band probe individual
axes, the Sobol decomposition then apportions the demand variance across all four
drivers at once, and the supply-side tests bound the cost axes the demand methods do not
reach. The conclusions below rest on all four methods together. The
reported bands rest on a 200-draw Monte Carlo per scenario. The convergence diagnostic
of Appendix~\ref{sec:method:uncertainty} finds the medians settled well inside that
sample and the deciles not converged at it, so the medians carry the results and the
80~per cent bands should be read as indicative of spread rather than as precise
intervals.

Two caveats bound the applicability of these methods and should be borne in mind
before the results are read. First, the per-NUTS3 building-stock feasibility
constraints on heat-pump and district-heat deployment are held deterministic;
probabilistic feasibility is reserved for future work. Second, the tornado is
one-at-a-time and does not capture covariation, so that, for example, a
simultaneous high carbon price and stranded hydrogen trajectory would reinforce and not offset one another. The Sobol screen apportions demand variance and
the tornado covers the three economic axes, but a full Morris screen of the
cost-side interactions is the identified next step.

\subsection{Tornado on the Three Economic Axes}
\label{sec:sensitivity:tornado}

Each of the three principal economic axes of the model, the ETS2 carbon-price
trajectory, the green-hydrogen import price and the discount rate, is moved between its
low and high bound while the other two are held at their central values. The screen is
deterministic. One central draw is used for every arm, with the same seed throughout, so
the only thing that differs between arms is the axis under test and the result is not
measuring sampling noise. The metric is buildings-sector carbon dioxide emissions in
2030 (MtCO$_2$ per year) under Stated Policies (Table~\ref{tab:tornado}). The near-term
horizon is deliberate, because by 2050 the cross-scenario spread is dominated by the
policy levers themselves.

The discount rate is the widest axis, at $-0.1$ to $+0.4$~MtCO$_2$, the ETS2 carbon
price next at $-0.3$ to $+0.1$, and the hydrogen price the narrowest at $-0.0$ to
$+0.3$. All three act through the levelised-cost layer and therefore only through the
cost-responsive part of the technology mix, which is what bounds them. The
variance-based global sensitivity of the demand side, which the tornado does not reach,
is set out in Appendix~\ref{sec:sensitivity:sobol} and
Figure~\ref{fig:sensitivity_panel}a.

\begin{table}[H]
    \centering
    \caption{Tornado on 2030 buildings-sector emissions (Stated~Policies), against a
    central-draw baseline of 539~MtCO$_2$. Each axis moves alone; every other lever is
    held and the model is re-run end to end. The screen is deterministic and is one draw,
    not a distribution.}
    \label{tab:tornado}
    \begin{tabular}{lcc}
        \toprule
        Axis & Emissions-lowering bound & Emissions-raising bound \\
        & (MtCO$_2$/yr) & (MtCO$_2$/yr) \\
        \midrule
        Discount rate     & $-0.1$ (low)        & $+0.4$ (high) \\
        ETS2 carbon price & $-0.3$ (high price) & $+0.1$ (low price) \\
        Hydrogen price    & $-0.0$ (rapid)      & $+0.3$ (stranded) \\
        \bottomrule
    \end{tabular}
    \\[2pt]{\footnotesize\textit{Source: Authors' contribution.}}
\end{table}

These magnitudes are small, and the reason is structural rather than a sign that the
price levers do not matter. The share allocation blends a scenario interpolation with a
cost-responsive term whose weight ramps linearly from zero in 2025 to 0.50 by 2050
(Appendix~\ref{sec:method:techshares}), so at 2030 only a tenth of the non-fossil
allocation responds to cost at all. A price axis moved alone at that date therefore
reaches very little of the mix, and the 2030 spread across the four scenarios, from 512
to 548~MtCO$_2$, is set by their share and demand levers rather than by the prices those
levers are paired with. The ordering the screen returns, with the discount rate widest and
hydrogen price narrowest, is itself informative. At 2030 the levelised comparison is
more sensitive to the cost of capital on equipment than to either fuel price, because
the capital charge is annualised over the whole equipment life while a fuel price moves
only the running cost of the small cost-responsive fraction.

The cost leadership behind these magnitudes is not uniform across countries, and it does
reverse in a few of them, so the tornado should be read as moving the size of a margin
that mostly runs one way, not as flipping the leader. The heat pump's advantage widens
in central and eastern Europe, where lower construction labour costs reduce heat-pump
capital expenditure and where cheaper retail electricity lowers its running cost, and
narrows in the cold, high-tariff north-west, where lower seasonal performance and dearer
electricity erode it. Against the incumbent gas boiler the best available heat pump is
cheaper in 24 of the 29 countries by 2050 on the central carbon path, losing only in
Cyprus, Croatia, Malta, Poland and the United Kingdom, and in 16 of 29 in 2025, so cost
leadership over fossil heating is broad but not universal at either date. That
comparison is the carbon-sensitive one, since the gas boiler pays ETS2 and the heat pump
does not, and its 2050 count runs from 19 of 29 on the low carbon path through 24 on the
central path to all 29 on the high path.

Against the hydrogen boiler the margin is narrower, and the 2050 levelised-cost ordering
of the heat pump, the hydrogen boiler and the gas boiler across all 29 countries is
reported in the Results section (Figure~\ref{fig:fig5}b). Which side of that ordering a
country falls on is governed by its storage geology first and by the hydrogen price path
second. Heat pumps are the cheaper base-load option in 13 of 29 countries on a rapid
hydrogen price decline, 22 on the central path, 26 on the slow path and all 29 if
hydrogen costs stall.

\subsection{Sweeping the Seasonal-Storage Adder}
\label{sec:sensitivity:sigma}

The same
storage adder decides the power arena's seven-market concentration, and it is
set by hand at \euro{}50/MWh-heat where salt caverns exist and \euro{}100 where they do
not, so an explanation of that concentration which appeals to the adder is circular
until the adder is itself varied. We vary it. Scaling the pair from 0.5 to 2.0 times
their central values, holding the cavern-to-non-cavern ratio at 2.0, moves the
power-arena win count as Table~\ref{tab:sigma_sweep} shows. The base-load count above is
not swept on this axis, so the base-load concentration is bounded here by the hydrogen
price paths alone.

\begin{table}[H]
\centering
\caption{Power-arena hydrogen wins under a swept seasonal-storage adder, H2~Push 2050.
The cavern and non-cavern adders are scaled together at a fixed ratio of 2.0. The count is invariant to the scarcity premium across \euro{}30 to 120/MWh-e, so one
column is shown. That invariance is an algebraic identity in the power arena, where the
premium enters both turbines' clearing price and cancels, rather than an empirical
result.}
\label{tab:sigma_sweep}
\small
\begin{tabular}{lrrrrrrrrr}
\toprule
Cavern adder (\euro{}/MWh-heat) & 25 & 38 & 45 & \textbf{50} & 55 & 62 & 75 & 88 & 100 \\
\midrule
Markets won, of 29              & 29 & 29 &  8 & \textbf{7}  &  7 &  7 &  7 &  2 &  0 \\
All winners cavern-endowed      & no & no & no & yes         & yes& yes& yes& yes& n/a \\
\bottomrule
\end{tabular}
\\[2pt]{\footnotesize\textit{Source: Authors' contribution.}}
\end{table}

The seven-market result is not an artefact of one point estimate, since it holds across
\euro{}50 to 75/MWh-heat, about a third of the \euro{}25 to 100 swept range. It is also
confined to that band. Below
\euro{}45 the count jumps to every market and the cavern confinement dissolves, and
above \euro{}75 it collapses towards zero. The concentration is a real property of the
model over a band; it is not a structural feature of the problem. One caveat bounds even that. The sweep scales the cavern
and non-cavern adders together at a fixed ratio, so the geological \emph{spread} that
actually drives the concentration is held constant throughout and is not tested here at all.

A second caveat bears on how the seven should be read. The classification is a two-value
national binary, and over the whole \euro{}50 to 75 band the winning set is that binary
exactly, so the seven is conditional on the classification rather than emerging
independently of it. Alternative readings of the same geological source move it. Adding
Iberia, which the source also maps, gives nine wins; restricting France to its northern
basins gives six. The defensible statement is therefore that, given a two-value national
classification of seven markets, the power-arena winners are exactly those seven, with a
range of six to nine under alternative readings. The ratio itself is corroborated
independently: a European techno-economic assessment of the same stores puts levelised
salt-cavern storage at about \$0.8/kg against \$1.5/kg for porous media
\citep{Talukdar2024UHSEurope}, a ratio near 1.9 against the 2.0 assumed here.

The carbon-price axis, by contrast, does not enter this particular
comparison at all, since neither a heat pump nor a hydrogen boiler pays a carbon price
on its own fuel, so the count is identical under low, central and high carbon
assumptions and the tornado's carbon axis acts only through the fossil technologies it
displaces. Regional labour-cost variation moves the size of the margin without setting
the ordering.

\subsection{Parametric Band on the Cross-Country Renovation Correlation ($\rho$)}
\label{sec:sensitivity:rho}

The Monte Carlo samples each country's envelope-renovation rate through a shared
EU-policy factor plus idiosyncratic national noise, with correlation $\rho$
(Appendix~\ref{sec:method:uncertainty}). The central 2050 demand trajectory is
invariant to $\rho$, because the shared factor preserves every country's marginal
distribution, so only the width of the band moves. At the central value
$\rho=0.5$, the EU 2050 figure is 3,651~TWh under Current Policies,
3,483~TWh under Stated Policies, 2,660~TWh under Net Zero and 3,154~TWh under H2
Push, essentially unchanged across $\rho$. The uncertainty band, by contrast,
widens
from 1.0$\times$ (independent, $\rho=0$) to 2.2$\times$ (shared, $\rho=0.5$) to
3.0$\times$ (fully correlated, $\rho=1$), by the same factor across all four
scenarios (Figure~\ref{fig:sensitivity_panel}b; the per-scenario bands are
collected in Table~\ref{tab:rho}). At $\rho=0$, independent national renovation
pace, the EU band is artificially narrow because national errors diversify away;
$\rho=0.5$, the shared central value, reflects a common policy environment under
the EPBD and the Renovation Wave. The demand bands should therefore be read as
conditional on $\rho=0.5$, while the headline scenario \emph{medians} are
unaffected.

\begin{table}[H]
    \centering
    \caption{2050 useful-heat demand band (TWh/yr, 10th to 90th percentile) against the
    cross-country renovation correlation $\rho$.}
    \label{tab:rho}
    \begin{tabular}{lccc}
        \toprule
        Scenario & $\rho=0$ (indep.) & $\rho=0.5$ (central) & $\rho=1$ (full) \\
        & (TWh/yr) & (TWh/yr) & (TWh/yr) \\
        \midrule
        Current Policies & 3,612 to 3,687 & 3,566 to 3,731 & 3,536 to 3,758 \\
        Stated Policies  & 3,432 to 3,533 & 3,369 to 3,593 & 3,328 to 3,630 \\
        Net Zero         & 2,563 to 2,770 & 2,435 to 2,900 & 2,353 to 2,977 \\
        H2 Push          & 3,080 to 3,229 & 2,987 to 3,319 & 2,928 to 3,375 \\
        \midrule
        Band-width vs $\rho=0$ & 1.0$\times$ & 2.2$\times$ & 3.0$\times$ \\
        \bottomrule
    \end{tabular}
    \\[2pt]{\footnotesize\textit{Source: Authors' contribution.}}
\end{table}

\begin{figure}[H]
    \centering
    \caption{Demand-side uncertainty structure. Panel (a) gives total-order ($S_T$, bars)
    and first-order ($S_1$, points) Sobol indices for 2050 useful heat; panel (b) gives the demand band width against the cross-country renovation correlation $\rho$, labelled at $\rho=1$.}
    \includegraphics[width=\textwidth]{paper/P15_sensitivity.pdf}
    \label{fig:sensitivity_panel}
    \\[2pt]{\footnotesize\textit{Source: Authors' contribution.}}
\end{figure}

\subsection{Global (Variance-Based) Sensitivity of 2050 Demand}
\label{sec:sensitivity:sobol}

The tornado and the parametric band above each move one axis at a time and so, by
construction, cannot apportion variance among the demand drivers or detect interactions
between them. The variance-based (Sobol) global sensitivity introduced here complements
them by decomposing the variance of EU useful heat in 2050 across all four uncertain
demand drivers at once: the envelope-renovation rate, the scenario lever, sampled over a
$\pm$35~per~cent band; population; occupancy, expressed as household size; and dwelling
size. The decomposition uses a Saltelli design with $N=2{,}048$. Because useful demand
in 2050 depends only on these drivers, it is evaluated analytically per country, so a
full Sobol design is both exact and immediate.

The total-order indices are an envelope rate of $S_T=0.40$, occupancy of $0.33$,
dwelling size of $0.16$ and population of $0.12$; the first-order indices approximately
equal the total-order indices, indicating a near-additive response with negligible
interactions. Two readings follow. First, the scenario lever, the renovation rate,
governs the largest single share of demand uncertainty, as intended. Second, occupancy,
the household-size trend that sets dwellings per capita, is almost as influential, which
is why the LMDI decomposition treats it as a first-class driver instead of folding it
into population. The cost and emissions axes, namely the carbon price, the hydrogen
trajectory, the discount rate, ETS2 pass-through and fuel-price correlation, act through
the levelised-cost layer, whose global screen is compute-bound on the full NUTS3 Monte
Carlo; the tornado above covers them, and a full Morris screen is the documented next
step.

\paragraph{Cost-side propagation.} The two caveats that bound this exercise, the
one-at-a-time tornado and the deterministic feasibility scores, are stated at the
head of the section. The cost-side axes are themselves
fully propagated. The hydrogen-price scenario, the discount rate, the carbon price,
ETS2 pass-through and the per-carrier log-normal fuel-price multipliers, covering
gas, electricity, oil and biomass, all enter the levelised-cost routine, which is
recomputed for each Monte Carlo sample, so the cost-side
uncertainty is active throughout even though its global variance decomposition,
which would require a screen of the full NUTS3 Monte Carlo, is left to the
identified Morris next step.

\subsection{Climate-Warming Demand Sensitivity}
\label{sec:sensitivity:climate}

The last of the demand-side methods bounds a structural assumption, not a sampled
parameter. The central demand trajectory holds heating degree days at the 1991 to 2020
climate normal, so it excludes the demand-reducing effect of a warming climate. Because space heating scales approximately linearly with heating degree days and domestic hot water does not, we overlay a warming factor on the space-heating share,
approximately 85 per cent of useful heat, using EU population-weighted degree-day
decline rates from \citet{Spinoni2018degreedays}: a central rate of $-0.35\%$ per year, which is that study's RCP4.5 result, swept over a $-0.20$ to $-0.55\%$ per year band. The upper end follows its RCP8.5 result; the study runs only those two pathways, so the lower end is our own conservative bound and not a reading of a lower-emissions scenario.

This is a one-sided robustness bound, not the central case. Under the central warming
rate, Stated Policies useful heat in 2050 falls a further 7~per~cent, from 3,483 to
3,235~TWh, deepening the 2025 to 2050 reduction from $-9\%$ to $-15.5\%$; relative to
the frozen-climate trajectory the additional 2050 reduction spans $-4.1\%$ at our own low
bound to $-10.9\%$ under RCP8.5. In the announced-policy pathway, a warming climate thus
delivers roughly half of the demand reduction that deep renovation delivers under Net
Zero ($-31\%$). The frozen-climate central trajectory is therefore conservative on
demand. The technology-mix and the RQ1 and RQ2 conclusions are unaffected, because
warming scales demand without altering relative levelised costs.

\subsection{Per-Country Hydrogen-Supply Cost and Import Connectivity}
\label{sec:sensitivity:h2supply}

The headline hydrogen prices apply a per-country multiplier to a single
north-west-European hub trajectory, which falls from \euro{}200/MWh in 2025 to
\euro{}50/MWh in 2050. The multiplier comes from
the bottom-up delivered-cost model of Appendix~\ref{sec:method:h2}, which costs out
five supply routes for each country and takes the cheapest, and normalises the result
so that the median country pays the hub price. This subsection reports what that model
finds.

Three findings emerge. First, pink and blue hydrogen are not the cheapest route in
2050. Nuclear electrolysis, at about \euro{}116 to 117/MWh and set by the cost of firm
nuclear power, and gas with carbon capture, at about \euro{}81/MWh
\citep{IEAGHG2022blue}, are both beaten by cheap green or imported hydrogen, so the
nuclear and carbon-capture `optionality' often invoked for central Europe does not in
fact lower its delivered cost. Second, domestic green electrolysis is the cheapest route
in the majority of the EU, with 23 countries self-supplying green hydrogen in 2050,
while the remaining six import by pipeline. No country takes the seaborne-ammonia
route at 2050, and the two non-backbone islands, Cyprus and Malta, self-supply green at
about \euro{}63 to 64/MWh against \euro{}108 for shipped ammonia
(Figure~\ref{fig:supplyrobust}a).

Third, capital and labour set the ranking, ahead of sunshine. High-WACC Greece, at
8~per~cent in the per-country cost-of-capital assumption of the hydrogen-supply
methodology (Appendix~\ref{sec:method:h2}), sits at a multiplier of 1.04 despite superb
solar resource, whereas cheap-capital, strong-wind Ireland, Denmark and the Baltic
states self-supply green hydrogen well below the median, at multipliers of 0.77 to 0.84.
Pipeline import caps the six connected countries at the hub price plus 9 to 16~per~cent, so
the dearest supply is found in the landlocked central-European states with mediocre
resource, namely Slovakia, the Czech Republic, Hungary, Slovenia and Italy at
multipliers of about 1.14 to 1.16, and in the non-backbone islands, Cyprus and Malta at
about 1.31 to 1.33, whose isolated-grid, small-scale solar is the costliest
self-supply of all and still undercuts the shipped-ammonia alternative.

Carrying each multiplier as a low, central and high band, the 2050
hydrogen-versus-heat-pump gap is negative, meaning hydrogen is cheaper, in a small
minority of countries. At central supply values, the default network treatment and
the seasonal store charged on the winter-weighted fraction of the hydrogen heating
load, 7 of the 29 countries sit at or below heat-pump parity
(Section~\ref{sec:results}). Cheap supply is not on its own enough to put a country
in that set. Finland and the Baltic states all sit below the median on the delivered
multiplier and none of them is in it, because none has the salt-cavern endowment
that makes seasonal storage cheap, which five of the seven that remain do have. Charging hydrogen
a fuller share of its distribution cost shrinks the set further
(Appendix~\ref{sec:results:h2infra}).

\paragraph{The colour of the molecule.} A second, qualitative finding runs with the
base-load result above, and it deserves more weight than its position in a sensitivity
section suggests. In nearly every assessed country the national hydrogen strategy
explicitly \emph{excludes} residential heating. More consequentially, the one dedicated assessment of hydrogen for buildings concludes that any heating hydrogen would in practice be \emph{blue} \citep{Dickel2024ET29}, and our own blue route prices that molecule at about \euro{}81/MWh \citep{IEAGHG2022blue}, against the green or imported molecule at about \euro{}50/MWh that the model prices. The model therefore
charges hydrogen a delivered commodity cost roughly \euro{}31/MWh below the only
sourcing the sectoral literature regards as realistic for this end use. That wedge is
wider than the widest of the seven surviving hydrogen leads, so a commodity price that
high would move every country further towards the heat pump.

On the storage-inclusive basis this paper uses, that repricing leaves the heat pump
cheaper in all 29 countries at a flat \euro{}81/MWh, and in 28 with Denmark
alone still holding if the German figure is read as a green-to-blue ratio rather than a
European price. The model's
delivered residential-hydrogen price is thus a counterfactual that it prices, not a
planned supply chain, and the counterfactual is generous to hydrogen. The base-load
result should be read in that light, as a statement about what a cheap green molecule
would do to the household economics of heating, not as a forecast that such a molecule
will be the one delivered to households.

\subsection{Electricity Distribution Reinforcement and the Scarcity Premium}
\label{sec:sensitivity:reinforce}

Two network-side assumptions deserve explicit stress tests. First, an accounting that
charges hydrogen for its distribution network only inside a dedicated
infrastructure scenario and never charges heat pumps for the low- and
medium-voltage grid \emph{reinforcement} that their added coincident winter peak
requires. That asymmetry has been removed. Both adders are now active by default in
every levelised cost this paper reports, so each carrier pays for its own last mile
(Appendix~\ref{sec:method:h2}). The reinforcement adder is
not opt-in; it is part of the headline number. The added
electric peak per dwelling, computed as annual heat over climate-scaled full-load
hours, divided by the cold-snap coefficient of performance, with a coincidence
factor of 0.5, is charged at \euro{}400, 800 and 1{,}600 per kW (low, central and
high), annualised at the country WACC over a 40-year network life. Across that band
the adder amounts to \euro{}1.1 to 38.2/MWh, and it is highest in peaky
mild-climate systems where a modest annual heat demand arrives through a short,
sharp season.

The stress test that follows is the conservative-for-heat-pump case, and it has two arms
that must be kept apart, because they charge hydrogen differently and return different
counts. Both charge the heat pump the \emph{high} reinforcement bound. In the first arm
hydrogen pays \emph{nothing} at all for its network, and heat pumps then remain cheaper
than hydrogen in 10 of 29 countries, namely Bulgaria, Estonia, Croatia, Hungary,
Lithuania, Latvia, Romania, Sweden, Slovenia and Slovakia, and cheaper than gas in 21 of
29, with a mean heat-pump-to-hydrogen margin of $-$\euro{}6.2/MWh and a median of
$-$\euro{}3.1/MWh, which is to say that hydrogen leads on average by those amounts. In
the second arm hydrogen pays the blended distribution network it is charged by default,
and heat pumps remain cheaper than hydrogen in 22 of 29 at a mean of $+$\euro{}13.6/MWh,
which reproduces the headline count at the top of the reinforcement band
(Figure~\ref{fig:supplyrobust}b).

The two arms should never be combined, since quoting the second count against the first
arm's network treatment would overstate the heat-pump case. The reading is the
same in either arm. The base-load ranking \emph{does} depend on how the two last-mile
networks are charged, and a reinforcement assumption at the top of its plausible band is
enough to move a majority of countries from one side of the comparison to the other.
This is one of the three most load-bearing assumptions in the levelised-cost layer,
alongside the seasonal-storage charge and the levy treatment of the two carriers, and
all three are carried explicitly in Appendix~\ref{sec:limitations}.

\begin{figure}[H]
    \centering
    \caption{Supply-side robustness in the hydrogen supply route and in the network
    charge. Panel~(a) gives each country's cheapest delivered hydrogen cost in 2050,
    coloured by which route wins, with the tick marking the next-best route so the margin
    between them is visible. Panel~(b) asks the question most favourable to hydrogen:
    hydrogen charged no last-mile network at all against the heat pump charged its high
    reinforcement band, per country, where a positive value still leaves the heat pump
    cheaper.}
    \includegraphics[width=\textwidth]{paper/P16_supply_robustness.pdf}
    \label{fig:supplyrobust}
    \\[2pt]{\footnotesize\textit{Source: Authors' contribution.}}
\end{figure}

Second, the merit-order results of Section~\ref{sec:results:meritorder} use an
endogenous scarcity electricity price, the marginal-peaker short-run marginal cost
plus a scarcity premium of \euro{}30, 60 and 120 per MWh (low, central and high).
The premium band spans the displacement of the thermal peaker by storage and
demand response across a varying number of hours. At the low premium the cold-snap heat pump becomes cheaper and hydrogen's
winning-country count at the building peak falls, while at the high premium it rises. The
counts are 0, 1, 8 and 15 of 29 across the four scenarios at \euro{}30/MWh-e, 0, 4, 9 and
16 at the central \euro{}60 and 0, 5, 10 and 18 at \euro{}120. Each sweep in this section
moves one lever and holds the rest at their published settings, so the baseline here is
the asymmetric accounting and the central column is its 0, 4, 9 and 16 rather than the
symmetric 0, 2, 7 and 14 the results section headlines; the last-mile charge is swept
separately. The demand-weighted
derived ceiling under Stated Policies runs from 0.2~per cent at the low premium to
5.3 at the high one against 5.1 at the central. That is the widest of the three
independent reasons the Stated Policies ceiling is a range rather than a point, and it is
the scenario in which the ceiling actually binds.
The power arena does not move at all, because there the premium enters both turbines'
clearing price and cancels. The missing-money conclusion is invariant across the band,
because it is the peaker's running hours, not the price level, that fail the
investment test.
